\chapter{Quantum Entanglements, Part 1}

These companion notes follow Leonard Susskind's fourth lecture in the \emph{Quantum Entanglement} supplementary track and are curated for this archive by LazyingArt LLC. The lecture opens by saying, in effect, that Bell's theorem will have to wait and that entanglement itself is only just within reach. But before two spins can do anything interesting together, we still have to finish the one-spin machinery. That is the governing decision of the lecture: first complete the story of a single two-state system, and only then return to the promised subject of entanglement.

\section{The Arbitrary-Axis Spin Operator}

We resume with a single spin and an arbitrary axis in space. Susskind labels the spatial directions by \(1,2,3\), while warning that he may lapse into \(x,y,z\). We will keep the \(1,2,3\) notation, with the understanding that these are just the three spatial axes.

An arbitrary direction is represented by a unit vector
\begin{equation}
\hat n=(n_1,n_2,n_3),
\qquad
n_1^2+n_2^2+n_3^2=1.
\end{equation}
If \(\sigma=(\sigma_1,\sigma_2,\sigma_3)\) denotes the spin operator, then the component of spin along \(\hat n\) is
\begin{equation}
\sigma\cdot\hat n = n_1\sigma_1+n_2\sigma_2+n_3\sigma_3.
\end{equation}
At this point the lecture makes an important distinction: this is the ordinary dot product in space, not the inner product of two state-vectors in Hilbert space. The \(n_i\) are just numbers, the direction cosines of the chosen axis. The \(\sigma_i\), however, are matrices. Therefore \(\sigma\cdot\hat n\) is itself an operator.

\begin{figure}[t]
\centering
\includegraphics[width=0.72\textwidth]{lecture_04_figure_02.png}
\caption{Spin component along an arbitrary direction.}
\end{figure}

Using the Pauli matrices, we obtain the explicit form
\begin{equation}
\sigma\cdot\hat n
=
\begin{pmatrix}
n_3 & n_1-i n_2 \\
n_1+i n_2 & -n_3
\end{pmatrix}.
\end{equation}
This is the operator that answers the physical question: if we orient the measuring apparatus along \(\hat n\), what are the possible results for the spin component along that axis?

\begin{figure}[t]
\centering
\includegraphics[width=0.78\textwidth]{lecture_04_figure_03.png}
\caption{Matrix form of \(\sigma\cdot\hat n\).}
\end{figure}

The lecture keeps the physical interpretation close at hand. To measure \(\sigma\cdot\hat n\), we imagine turning on a large magnetic field oriented along \(\hat n\). The electron either emits a photon or does not. That physical process is encoded mathematically by the operator \(\sigma\cdot\hat n\). The whole point is that the abstract operator language and the laboratory procedure are describing the same thing.

\section{Basic Properties of \texorpdfstring{$\sigma\cdot\hat n$}{sigma dot n}}

It is convenient to introduce the standard shorthand
\begin{equation}
n_- = n_1-i n_2,
\qquad
n_+ = n_1+i n_2.
\end{equation}
Then the matrix can be written as
\begin{equation}
\sigma\cdot\hat n
=
\begin{pmatrix}
n_3 & n_- \\
n_+ & -n_3
\end{pmatrix}.
\end{equation}
The lower shorthand line on the board is only partially visible, but the notation is clear from the surrounding lecture.

The matrix is Hermitian: the diagonal entries are real, and the off-diagonal entries are complex conjugates of one another. That is exactly what we want from an observable. The lecture then extracts a few working facts and moves on; these are tools, not a separate formalism.

\begin{remark}
Here and below, the symbol \(1\) means the identity operator on the two-dimensional one-spin state space.
\end{remark}

\begin{proposition}
For any unit vector \(\hat n\),
\begin{equation}
(\sigma\cdot\hat n)^2=1.
\end{equation}
Hence the eigenvalues of \(\sigma\cdot\hat n\) are \(\pm1\).
\end{proposition}

\begin{proof}
One may verify the square directly from the matrix and the identity \(n_1^2+n_2^2+n_3^2=1\). The lecture also recalls the Pauli-matrix relation
\begin{equation}
\sigma_i\sigma_j=-\sigma_j\sigma_i,
\qquad i\neq j,
\end{equation}
which gives the same result. If \(\lambda\) is an eigenvalue of \(\sigma\cdot\hat n\), then \(\lambda^2=1\), so \(\lambda=\pm1\).
\end{proof}

A second algebraic identity is introduced because it will be used repeatedly in the eigenvector calculation:
\begin{equation}
n_+n_-=(n_1+i n_2)(n_1-i n_2)=n_1^2+n_2^2=1-n_3^2.
\end{equation}
The lecture treats this exactly as it should be treated: not as a deep theorem, but as the small fact that makes the later algebra survivable.

\section{From Preparation to Probability}

Now the lecture arrives at the question it really wants to answer. Suppose we prepare the electron so that its spin is definitely up along the axis \(\hat n\). Then we measure the spin along a different axis \(\hat m\). What is the probability that the outcome along \(\hat m\) is \(+1\)?

Before calculating anything, Susskind notes what symmetry suggests. The answer ought to depend only on the angle between \(\hat n\) and \(\hat m\), and it ought to be symmetric in \(n\) and \(m\). But he does not stop there. The lecture insists that we derive the answer from the probability postulate rather than simply guess the angle-dependence and bless it afterward.

We therefore introduce the \(+1\) eigenstates of the two operators:
\begin{equation}
\sigma\cdot\hat n\,|\sigma\cdot\hat n=1\rangle
=
|\sigma\cdot\hat n=1\rangle,
\qquad
\sigma\cdot\hat m\,|\sigma\cdot\hat m=1\rangle
=
|\sigma\cdot\hat m=1\rangle.
\end{equation}
The equality sign inside each ket is part of the label; it is not an additional assertion. The lecture is explicit about this because the notation is a bit heavy. The board is partially occluded at this point, so what we write here is the cautious standard reconstruction of the intended notation.

\begin{figure}[t]
\centering
\includegraphics[width=0.72\textwidth]{lecture_04_figure_04.png}
\caption{Spin-up eigenstates along \(\hat m\) and \(\hat n\).}
\end{figure}

\begin{figure}[t]
\centering
\begin{tikzpicture}[scale=1.4]
\draw[->] (0,0) -- (1.3,0.25) node[right] {$\hat n$};
\draw[->] (0,0) -- (0.65,1.18) node[above] {$\hat m$};
\end{tikzpicture}
\caption{A clean reconstruction of the \(\hat n\)-\(\hat m\) sketch that remains on the board while the eigenstate equations are written.}
\end{figure}

\subsection{Question \& Answer}

\paragraph{Question.} Once the state is prepared along \(\hat n\), what quantity do we calculate first in order to predict the outcome along \(\hat m\)?

\paragraph{Answer.} First we calculate a probability amplitude, not a probability. The amplitude is the inner product between the prepared state and the final eigenstate whose outcome we want to test:
\begin{equation}
\langle \sigma\cdot\hat m=1 \mid \sigma\cdot\hat n=1\rangle.
\end{equation}
Only after that do we form the probability by multiplying by the complex conjugate:
\begin{equation}
P(\sigma\cdot\hat m=1 \mid \sigma\cdot\hat n=1)
=
\left|\langle \sigma\cdot\hat m=1 \mid \sigma\cdot\hat n=1\rangle\right|^2.
\end{equation}
That distinction is one of the lecture's central pedagogical beats: amplitude first, probability second.

\begin{figure}[t]
\centering
\includegraphics[width=0.74\textwidth]{lecture_04_figure_05.png}
\caption{Overlap of spin-up states along \(\hat m\) and \(\hat n\).}
\end{figure}

The transcript is garbled around this point, but the board progression and the surrounding discussion make the intended route clear: define the relevant eigenvectors, take their overlap, then square the magnitude.

\section{The Eigenvector Calculation and the Formula \texorpdfstring{$\frac{1+\hat n\cdot\hat m}{2}$}{(1+n.m)/2}}

The lecture now slows down and performs the calculation. Susskind warns the audience that it is ugly. We may clean it up a little, but we should preserve the rhythm of the argument: first solve for the eigenvectors, then normalize them, and only then take the inner product.

We seek a \(+1\) eigenvector of \(\sigma\cdot\hat n\), so we solve
\begin{equation}
(\sigma\cdot\hat n)\binom{\alpha}{\beta}=\binom{\alpha}{\beta}.
\end{equation}
At first we do \emph{not} normalize the vector. We also keep in mind another point the lecture emphasizes: multiplying a state by an overall phase \(e^{i\theta}\) changes no probability, since every physical probability is built from a quantity times its complex conjugate. Therefore, at this stage, there is no loss of physical content in choosing
\begin{equation}
\alpha=1.
\end{equation}
Write the lower component as \(z\). Then
\begin{equation}
\begin{pmatrix}
n_3 & n_- \\
n_+ & -n_3
\end{pmatrix}
\begin{pmatrix}
1\\
z
\end{pmatrix}
=
\begin{pmatrix}
1\\
z
\end{pmatrix}.
\end{equation}
From the top component we obtain
\begin{equation}
n_3+n_-z=1,
\end{equation}
and hence
\begin{equation}
z=\frac{1-n_3}{n_-}.
\end{equation}
So one unnormalized \(+1\) eigenvector is
\begin{equation}
\begin{pmatrix}
1\\[4pt]
\dfrac{1-n_3}{n_-}
\end{pmatrix}.
\end{equation}

Now we normalize it. The corresponding bra is
\begin{equation}
\left(
1,\frac{1-n_3}{n_+}
\right),
\end{equation}
so the norm is
\begin{align}
\left(
1,\frac{1-n_3}{n_+}
\right)
\begin{pmatrix}
1\\[4pt]
\dfrac{1-n_3}{n_-}
\end{pmatrix}
&=
1+\frac{(1-n_3)^2}{n_+n_-} \\
&=
1+\frac{(1-n_3)^2}{(1-n_3)(1+n_3)} \\
&=
\frac{2}{1+n_3}.
\end{align}
Therefore a convenient normalized eigenvector is
\begin{equation}
|\sigma\cdot\hat n=1\rangle
=
\sqrt{\frac{1+n_3}{2}}
\begin{pmatrix}
1\\[6pt]
\dfrac{1-n_3}{n_-}
\end{pmatrix}.
\end{equation}
This is one explicit phase convention. The lecture would not want us to forget that other equivalent conventions exist; what matters physically is the ray, not the overall phase.

The corresponding bra for the \(+1\) eigenstate of \(\sigma\cdot\hat m\) is obtained by replacing \(n\) with \(m\) and taking the conjugate:
\begin{equation}
\langle \sigma\cdot\hat m=1|
=
\left(
\sqrt{\frac{1+m_3}{2}},
\ 
\frac{1-m_3}{m_+}\sqrt{\frac{1+m_3}{2}}
\right).
\end{equation}

The next step is the one Susskind refuses to do line by line on the board. We form the inner product
\begin{equation}
\langle \sigma\cdot\hat m=1 \mid \sigma\cdot\hat n=1\rangle
\end{equation}
and then multiply by its complex conjugate. The intermediate expression is, in his words, an unholy mess. But after repeated use of identities such as \(n_+n_-=1-n_3^2\), everything collapses to
\begin{equation}
P(\sigma\cdot\hat m=1 \mid \sigma\cdot\hat n=1)
=
\frac{1+n_1m_1+n_2m_2+n_3m_3}{2}
=
\frac{1+\hat n\cdot\hat m}{2}.
\end{equation}
If \(\theta_{nm}\) is the angle between the two axes, then
\begin{equation}
\hat n\cdot\hat m=\cos\theta_{nm},
\end{equation}
and therefore
\begin{equation}
P(\sigma\cdot\hat m=1 \mid \sigma\cdot\hat n=1)
=
\frac{1+\cos\theta_{nm}}{2}.
\end{equation}

Now the lecture does what it always does after a calculation: it checks the answer against obvious cases.
\begin{itemize}
\item If \(\hat m=\hat n\), then \(\cos\theta_{nm}=1\), so \(P=1\).
\item If \(\hat m=-\hat n\), then \(\cos\theta_{nm}=-1\), so \(P=0\).
\item If \(\hat m\perp\hat n\), then \(\cos\theta_{nm}=0\), so \(P=\tfrac12\).
\end{itemize}
So the formula passes the tests the lecture wanted: same axis gives certainty, opposite axis gives impossibility, and a right angle gives a half-half result.

\begin{remark}
After the fact, one can simplify the derivation by choosing coordinates so that \(\hat n\) is the third axis. Then the prepared state is simply \(\binom{1}{0}\), and the amplitude is immediately the upper component of the \(m\)-eigenvector. The lecture deliberately resists doing that at the start, because the point is to \emph{derive} the rotationally invariant answer rather than assume it.
\end{remark}

The lecture also pauses to interpret what probability means operationally. One does not measure a probability in a single shot. One either repeats the same preparation-and-measurement sequence on the same electron, re-preparing it each time along \(\hat n\), or else prepares many identical electrons and performs the same measurement on each. In either method the fraction of \(+1\) outcomes tends to the probability above.

The real surprise is now in plain view. Every measurement of spin along any chosen axis still gives only \(\pm1\), never anything in between. Yet the theory is rotationally smooth, not because the individual outcomes are fractional, but because the \emph{probabilities} are.

\section{Measurement as Filter; Simultaneous Measurements}

With the probability formula in hand, the lecture turns to another important point that had been hovering in the background: a measurement changes the state.

If we measure \(\sigma\cdot\hat m\) and the outcome is \(+1\), we do not merely discover something about the original state. We also leave the electron in the state \(|\sigma\cdot\hat m=1\rangle\). The lecture compares this to a polarizer for photons: the device both tests the polarization and prepares the outgoing polarization state. A measurement is therefore a filter.

This is the bridge to simultaneous measurement. If one measurement leaves the system in an eigenstate of the measured operator, then a second quantity can be measured simultaneously only if the same post-measurement state is an eigenstate of both observables.

\subsection{Question \& Answer}

\paragraph{Question.} When can two observables be measured simultaneously?

\paragraph{Answer.} They must have the same eigenvectors. Otherwise the state left behind by one measurement cannot also be a definite state for the other observable.

This is why different spin components do not admit simultaneous reproducible measurement. For \(\sigma\cdot\hat n\) and \(\sigma\cdot\hat m\) with \(\hat n\neq \hat m\), the eigenvectors are different. There is no state that is simultaneously an eigenvector of both. So one cannot measure both without spoiling the information carried by the other.

The lecture then writes the general theorem on the board.

\begin{figure}[t]
\centering
\includegraphics[width=0.72\textwidth]{lecture_04_figure_06.png}
\caption{Two matrices sharing the same eigenvectors.}
\end{figure}

\begin{theorem}
Suppose \(A\) and \(B\) share the same complete set of eigenvectors \(|\alpha_i\rangle\):
\begin{equation}
A|\alpha_i\rangle=\lambda_i|\alpha_i\rangle,
\qquad
B|\alpha_i\rangle=\mu_i|\alpha_i\rangle.
\end{equation}
Then
\begin{equation}
AB=BA.
\end{equation}
Equivalently,
\begin{equation}
[A,B]\equiv AB-BA=0.
\end{equation}
\end{theorem}

\begin{proof}
On any common eigenvector,
\begin{equation}
AB|\alpha_i\rangle
=
A(\mu_i|\alpha_i\rangle)
=
\lambda_i\mu_i|\alpha_i\rangle,
\end{equation}
while
\begin{equation}
BA|\alpha_i\rangle
=
B(\lambda_i|\alpha_i\rangle)
=
\mu_i\lambda_i|\alpha_i\rangle.
\end{equation}
Thus \(AB\) and \(BA\) act the same way on every basis vector \(|\alpha_i\rangle\), so \(AB=BA\).
\end{proof}

The lecture states the converse as well in the present finite-dimensional setting: commuting operators can be simultaneously diagonalized. That is the standard criterion for simultaneous measurability. The proof of the converse is not developed in detail here, but the theorem itself is part of the lecture's mathematical payload.

\section{Any One-Spin State Is Up Along Some Axis}

There is one more single-spin theorem before the lecture will allow itself to move on to entanglement. Up to now we have been given a direction \(\hat n\) and have found the corresponding eigenvector. The lecture now reverses the problem.

Take an arbitrary normalized spinor
\begin{equation}
\begin{pmatrix}
\alpha\\
\beta
\end{pmatrix},
\qquad
|\alpha|^2+|\beta|^2=1.
\end{equation}
Can we find a direction \(\hat n\) such that
\begin{equation}
(\sigma\cdot\hat n)
\begin{pmatrix}
\alpha\\
\beta
\end{pmatrix}
=
\begin{pmatrix}
\alpha\\
\beta
\end{pmatrix}?
\end{equation}

\subsection{Question \& Answer}

\paragraph{Question.} Given an arbitrary normalized one-spin state, is there always some axis along which the spin is definitely up?

\paragraph{Answer.} Yes. The lecture proves it by counting degrees of freedom.

A unit vector \(\hat n\) has only two independent real parameters, since \(n_1^2+n_2^2+n_3^2=1\). On the other hand, the state-vector \(\binom{\alpha}{\beta}\) begins with four real parameters, but normalization removes one and overall phase removes one more. Thus a physical one-spin state also depends on only two real parameters.

So the reverse eigenvalue problem has the right number of meaningful unknowns. In the lecture's language, there are two equations and two variables. Therefore there is always some direction \(\hat n\) for which the given spinor is the \(+1\) eigenstate of \(\sigma\cdot\hat n\).

This is a concrete completion of the one-spin story. No matter how a pure one-spin state is prepared, it can always be thought of as ``spin up'' along some axis.

The lecture even gives an operational picture. Prepare a large number of identical copies of the same electron state. Measure one batch along one axis, another batch along a nearby axis, and keep adjusting the direction of the magnet until you find the axis along which all test electrons come out up. That is the axis corresponding to the state.

At this point the one-spin machinery is complete: arbitrary-axis observables, probability amplitudes, the measurement-as-filter rule, the commutator criterion, and the reverse theorem that identifies every pure one-spin state with spin-up along some direction. Only now does the lecture return to the topic it promised at the start.

\section{Two Spins, Product States, and the First Entangled State}

We now begin the two-spin system. The lecture simplifies the situation brutally: forget the motion of the electrons; imagine them nailed down in space so that only their spins are free to vary. That is enough to reach entanglement.

Just before switching to two spins, Susskind writes a compact action table for the Pauli matrices on the \(\sigma_3\)-basis states \(|u\rangle\) and \(|d\rangle\):
\begin{align}
\sigma_1|u\rangle &= |d\rangle, &
\sigma_1|d\rangle &= |u\rangle, \\
\sigma_2|u\rangle &= i|d\rangle, &
\sigma_2|d\rangle &= -i|u\rangle, \\
\sigma_3|u\rangle &= |u\rangle, &
\sigma_3|d\rangle &= -|d\rangle.
\end{align}
This is the same content as the matrix representation, now repackaged in the language of kets. The scalar factors \(i\) and \(-i\) multiply the vector from outside; they do not alter the label inside the ket.

For two electrons there are four basis states:
\begin{equation}
|uu\rangle,\qquad |du\rangle,\qquad |ud\rangle,\qquad |dd\rangle.
\end{equation}
The first slot refers to electron \(1\), the second to electron \(2\). These four states correspond to distinct measurement outcomes of the third components of the two spins, and the lecture treats them as mutually orthogonal basis vectors.

The first electron's spin operators continue to be called \(\sigma_1,\sigma_2,\sigma_3\). The second electron's are renamed \(\tau_1,\tau_2,\tau_3\). Their action is simple:
\begin{itemize}
\item \(\sigma_i\) acts only on the first slot and ignores the second;
\item \(\tau_i\) acts only on the second slot and ignores the first.
\end{itemize}
For example,
\begin{equation}
\sigma_1|ud\rangle=|dd\rangle,
\qquad
\sigma_1|dd\rangle=|ud\rangle,
\qquad
\sigma_3|dd\rangle=-|dd\rangle,
\qquad
\tau_2|ud\rangle=-i|uu\rangle.
\end{equation}
That is the lecture's rule: the untouched electron simply goes along for the ride.

\subsection{Question \& Answer}

\paragraph{Question.} Why does the two-spin system require six real parameters rather than four?

\paragraph{Answer.} Because the general two-spin state is richer than ``electron one in its own state'' plus ``electron two in its own state.''

A general state of two spins is
\begin{equation}
|\Psi\rangle
=
c_{uu}|uu\rangle
+
c_{du}|du\rangle
+
c_{ud}|ud\rangle
+
c_{dd}|dd\rangle,
\end{equation}
with four complex coefficients. That is eight real numbers. Normalization removes one real parameter and overall phase removes another, leaving
\begin{equation}
8-1-1=6
\end{equation}
independent real parameters.

If we had merely added the two physical parameters of the first electron to the two physical parameters of the second, we would have guessed only \(4\). The fact that the actual answer is \(6\) is exactly the pressure that drives the lecture to entanglement: there are more states than product states.

Now we describe the product states first, because the lecture wants the contrast to be sharp. If electron \(1\) is in the state
\begin{equation}
\alpha_1|u\rangle+\beta_1|d\rangle,
\end{equation}
and electron \(2\) is in the state
\begin{equation}
\alpha_2|u\rangle+\beta_2|d\rangle,
\end{equation}
then the combined two-electron state is their product:
\begin{equation}
(\alpha_1|u\rangle+\beta_1|d\rangle)(\alpha_2|u\rangle+\beta_2|d\rangle).
\end{equation}
Expanding it in the two-spin basis gives
\begin{align}
(\alpha_1|u\rangle+\beta_1|d\rangle)(\alpha_2|u\rangle+\beta_2|d\rangle)
&=
\alpha_1\alpha_2|uu\rangle
+
\alpha_1\beta_2|ud\rangle \\
&\quad
+
\beta_1\alpha_2|du\rangle
+
\beta_1\beta_2|dd\rangle.
\end{align}
Such a state is a \emph{product state}. Each electron still has its own definite polarization axis, just as in the one-spin story. Product states are the natural states of two independently prepared electrons.

The lecture then constructs a state which manifestly does \emph{not} have that property:
\begin{equation}
|\psi_+\rangle
=
\frac{1}{\sqrt2}\bigl(|ud\rangle+|du\rangle\bigr).
\end{equation}
This is the first entangled example in the lecture, and it is important to keep the plus sign here. Only later does Susskind remark that the minus-sign version is the better state and will matter next time.

The first property of this state is immediate. If we measure electron \(2\) and find it down, then electron \(1\) must be up. If we find electron \(2\) up, then electron \(1\) must be down. So the state carries definite anti-correlation in the \(3\)-direction.

But the lecture does not stop with that observation. It proves something more startling: every one-spin expectation value vanishes in this state.

For \(\sigma_1\),
\begin{align}
\sigma_1|\psi_+\rangle
&=
\frac{1}{\sqrt2}\bigl(|dd\rangle+|uu\rangle\bigr), \\
\langle \psi_+|\sigma_1|\psi_+\rangle
&=
\frac{1}{2}
(\langle ud|+\langle du|)(|dd\rangle+|uu\rangle)
=0.
\end{align}
No term survives because the basis states are orthogonal: the bra contains only \(\langle ud|\) and \(\langle du|\), while the ket contains only \(|dd\rangle\) and \(|uu\rangle\).

The same orthogonality argument gives
\begin{equation}
\langle \psi_+|\sigma_2|\psi_+\rangle=0.
\end{equation}
For \(\sigma_3\), the cancellation is slightly different:
\begin{align}
\sigma_3|\psi_+\rangle
&=
\frac{1}{\sqrt2}\bigl(|ud\rangle-|du\rangle\bigr), \\
\langle \psi_+|\sigma_3|\psi_+\rangle
&=
\frac{1}{2}
(\langle ud|+\langle du|)(|ud\rangle-|du\rangle)
=0.
\end{align}
Thus
\begin{equation}
\langle \psi_+|\sigma_i|\psi_+\rangle=0,
\qquad i=1,2,3.
\end{equation}
By exactly the same reasoning,
\begin{equation}
\langle \psi_+|\tau_i|\psi_+\rangle=0,
\qquad i=1,2,3.
\end{equation}
Hence for any direction \(\hat n\),
\begin{equation}
\langle \psi_+|\sigma\cdot\hat n|\psi_+\rangle
=
n_1\langle \sigma_1\rangle
+
n_2\langle \sigma_2\rangle
+
n_3\langle \sigma_3\rangle
=
0,
\end{equation}
and likewise for the second electron.

This is the genuinely new phenomenon. For a single pure spin state, there is always some axis along which the spin is definitely up. But for this two-electron state, every one-spin direction gives a half-half distribution. No axis is preferred for either electron by itself, even though the pair carries definite correlation information.

The lecture now asks how odd this really is. At first sight, not very. If we mix up a penny and a dime, separate them without looking, and inspect only one, then the result tells us what the other must be. So correlation by itself is not yet the deepest mystery. The penny-dime analogy is used very carefully: it shows that simple anti-correlation is not enough to distinguish the classical from the quantum. The real difference is Bell's theorem, and the lecture explicitly postpones that difference to the next time.

There is one final turn. Near the end, Susskind says that the minus-sign combination
\begin{equation}
\frac{1}{\sqrt2}\bigl(|ud\rangle-|du\rangle\bigr)
\end{equation}
is the better state. He sketches a physical reason why two nearby electron magnets tend to anti-align and even names the later state as the singlet. But the reason for the minus sign is also declared to be non-obvious and is deferred. So the present chapter should end exactly where the lecture ends: with the first entangled state understood, with the better state visible on the horizon, and with Bell still waiting.

\section{Summary}

The lecture began by promising entanglement and then refusing to rush to it. That refusal was mathematically productive. We first completed the one-spin machinery: the operator \(\sigma\cdot\hat n\), its explicit matrix form, its eigenvalues and eigenvectors, the distinction between amplitude and probability, and the formula
\begin{equation}
P(\sigma\cdot\hat m=1 \mid \sigma\cdot\hat n=1)=\frac{1+\hat n\cdot\hat m}{2}.
\end{equation}
Along the way we clarified two general principles: a measurement leaves the system in an eigenstate, and simultaneous measurability is tied to commuting operators.

Only after that could two spins reveal something new. The two-electron state space is richer than the space of simple product states, and the symmetric state \(\frac{1}{\sqrt2}(|ud\rangle+|du\rangle)\) shows the first unmistakable sign of that richness. Each electron by itself looks completely unbiased in every direction, yet measuring one teaches us something about the other. That is the lecture's first real encounter with quantum entanglement, and it is left poised just before Bell's theorem and just before the full importance of the minus-sign singlet state.